\documentclass[10pt]{article}

\usepackage{xr-hyper}
\externaldocument[][nocite]{main}
\usepackage[T1]{fontenc}
\usepackage[utf8]{inputenc}
\usepackage{lmodern}
\usepackage{microtype}
\usepackage[a4paper,margin=1in]{geometry}

\usepackage{amsmath,amssymb,amsthm,mathtools}
\usepackage{bm}
\usepackage{booktabs}
\usepackage{array}
\usepackage{tabularx}
\usepackage{longtable}
\usepackage{enumitem}
\usepackage{xcolor}
\usepackage{graphicx}
\usepackage{url}
\usepackage{natbib}
\usepackage{hyperref}
\usepackage[nameinlink,capitalise,noabbrev]{cleveref}

\definecolor{deepblue}{HTML}{173B57}
\definecolor{softblue}{HTML}{EAF3F8}
\definecolor{deepgreen}{HTML}{235D3A}
\definecolor{softgreen}{HTML}{EAF5EE}
\definecolor{deeporange}{HTML}{9C4E00}
\definecolor{softorange}{HTML}{FFF1DF}
\definecolor{draftred}{HTML}{8F1D1D}
\definecolor{draftbg}{HTML}{FFF4F4}
\definecolor{openpurple}{HTML}{5B3F8C}
\definecolor{softpurple}{HTML}{F2ECFA}
\definecolor{graytext}{HTML}{4B5563}

\hypersetup{
  colorlinks=true,
  linkcolor=deepblue,
  citecolor=deepgreen,
  urlcolor=deepblue,
  pdftitle={Nuclear-Norm Best Approximation at Rank-Deficient Matrix-Line Crossings},
  pdfauthor={Anonymous manuscript}
}

\setlength{\parindent}{0pt}
\setlength{\parskip}{0.55em}
\setlist[itemize]{topsep=0.35em,itemsep=0.2em}
\setlist[enumerate]{topsep=0.35em,itemsep=0.25em}
\renewcommand{\arraystretch}{1.15}

\newtheoremstyle{reader}
  {0.8em}{0.8em}
  {\itshape}{}
  {\bfseries\color{deepblue}}{.}
  {0.5em}{}
\theoremstyle{reader}
\newtheorem{theorem}{Theorem}[section]
\newtheorem{proposition}[theorem]{Proposition}
\newtheorem{lemma}[theorem]{Lemma}
\newtheorem{corollary}[theorem]{Corollary}

\theoremstyle{definition}
\newtheorem{definition}[theorem]{Definition}
\newtheorem{assumption}[theorem]{Assumption}
\newtheorem{example}[theorem]{Example}

\theoremstyle{remark}
\newtheorem{remark}[theorem]{Remark}

\newenvironment{proofidea}
  {\par\smallskip\begin{center}\begin{minipage}{0.94\textwidth}
   {\color{deepblue}\hrule}\vspace{0.35em}
   \textbf{\color{deepblue}Proof idea.}\ \color{black}\ignorespaces}
  {\par\vspace{0.35em}{\color{deepblue}\hrule}
   \end{minipage}\end{center}\smallskip}

\newenvironment{intuition}
  {\par\smallskip\begin{center}\begin{minipage}{0.94\textwidth}
   {\color{deepgreen}\hrule}\vspace{0.35em}
   \textbf{\color{deepgreen}What this result does.}\ \color{black}\ignorespaces}
  {\par\vspace{0.35em}{\color{deepgreen}\hrule}
   \end{minipage}\end{center}\smallskip}

\newenvironment{boundary}
  {\par\smallskip\begin{center}\begin{minipage}{0.94\textwidth}
   {\color{deeporange}\hrule}\vspace{0.35em}
   \textbf{\color{deeporange}Boundary of the claim.}\ \color{black}\ignorespaces}
  {\par\vspace{0.35em}{\color{deeporange}\hrule}
   \end{minipage}\end{center}\smallskip}

\newenvironment{openproblem}
  {\par\smallskip\begin{center}\begin{minipage}{0.94\textwidth}
   {\color{openpurple}\hrule}\vspace{0.35em}
   \textbf{\color{openpurple}Open problem.}\ \color{black}\ignorespaces}
  {\par\vspace{0.35em}{\color{openpurple}\hrule}
   \end{minipage}\end{center}\smallskip}

\newcommand{\proved}{\textcolor{deepgreen}{\textbf{PROVED}}}
\newcommand{\localproved}{\textcolor{deeporange}{\textbf{PROVED UNDER LOCAL ASSUMPTIONS}}}
\newcommand{\observed}{\textcolor{deepblue}{\textbf{NUMERICALLY OBSERVED}}}
\newcommand{\openstatus}{\textcolor{openpurple}{\textbf{OPEN}}}

\newcommand{\ip}[2]{\left\langle #1,#2\right\rangle}
\newcommand{\R}{\mathbb{R}}
\newcommand{\norm}[1]{\left\lVert #1\right\rVert}
\newcommand{\snorm}[1]{\left\lVert #1\right\rVert_{2}}
\newcommand{\fnorm}[1]{\left\lVert #1\right\rVert_{\mathrm F}}
\newcommand{\nnorm}[1]{\left\lVert #1\right\rVert_{*}}
\newcommand{\argmin}{\operatorname*{arg\,min}}
\newcommand{\argmax}{\operatorname*{arg\,max}}
\newcommand{\diag}{\operatorname{diag}}
\newcommand{\rank}{\operatorname{rank}}
\newcommand{\range}{\operatorname{range}}
\newcommand{\Retr}{\operatorname{Retr}}
\newcommand{\Gap}{\operatorname{Gap}}
\newcommand{\Sc}{S_{\mathrm c}}
\newcommand{\Kcan}{K_{\mathrm{can}}}
\newcommand{\Pstar}{\mathcal P^\star}
\newcommand{\Ctheta}{\mathcal C_{\Theta}}

\DeclareMathOperator{\tr}{tr}
\DeclareMathOperator{\dist}{dist}
\DeclareMathOperator{\sign}{sign}

\crefname{theorem}{theorem}{theorems}
\crefname{proposition}{proposition}{propositions}
\crefname{lemma}{lemma}{lemmas}
\crefname{corollary}{corollary}{corollaries}
\crefname{definition}{definition}{definitions}
\crefname{assumption}{assumption}{assumptions}
\crefname{example}{example}{examples}
\crefname{remark}{remark}{remarks}

\hypersetup{
  pdftitle={Technical Supplement: Nuclear-Norm Best Approximation at Rank-Deficient Matrix-Line Crossings},
  pdfauthor={Anonymous manuscript}
}
\setlength{\parskip}{0.35em}
\renewcommand{\thesection}{S\arabic{section}}

\title{\textbf{Technical Supplement to}\
\textbf{Nuclear-Norm Best Approximation at Rank-Deficient Matrix-Line Crossings}}
\author{Anonymous manuscript}
\date{August 4, 2026}

\begin{document}
\maketitle

\section{Index and scope of the supplement}

This supplement contains auxiliary results that are useful for implementing
or embedding the inner oracle but are not needed for the main theorem chain.
The essential proofs of compact-root failure, certificate completion, Fenchel
recovery/selection, and the singular rates remain in the main paper.

\begin{enumerate}
 \item \textbf{Exact interval bisection.}  We record the elementary
 convergence proof for the set-valued scalar update used by the exact
 reference oracle.
 \item \textbf{Spectral-level-set normal geometry.}  We derive the simple-top
 normal, distinguish the repeated-top Bouligand tangent cone, quantify radial
 drift, and give exact projective normal-error formulas.  This material
 justifies the rank-one normal in the motivating
 spectral-level-set application; the main matrix-line results permit general
 normals.
 \item \textbf{Feasible primal--dual certificate.}  We give an exact-real
 repair and a gap decomposition for an approximate smooth direction, while
 separating those identities from the stronger requirements of a
 machine-validated implementation.
 \item \textbf{Detailed reproducibility tables.}  We report the full
 constant-level checks for the regular, strict, boundary, moving-crossover,
 smoothing-family, direction, and higher-order cancellation formulas
 summarized in the main paper.
\end{enumerate}

\section{Exact interval bisection}

Let \(L(\lambda)=a(\lambda)-\beta(\lambda)\) and
\(U(\lambda)=a(\lambda)+\beta(\lambda)\), with the notation of
Theorem~\ref{thm:interval}.

\begin{proposition}[Set-valued bisection]
If \(U(\lambda)<0\), every minimizer lies to the right of \(\lambda\); if
\(L(\lambda)>0\), every minimizer lies to the left; and if
\(L(\lambda)\le0\le U(\lambda)\), then \(\lambda\) is a minimizer.  Starting
from the bracket in \eqref{eq:lambda-bound}, midpoint updates either stop at
an exact minimizer or generate nested intervals converging to one.
\end{proposition}

\begin{proof}
The first two cases say that the entire scalar subdifferential is strictly
negative or strictly positive, so one-dimensional convexity identifies the
side containing every minimizer.  The third case is
\(0\in\partial\phi(\lambda)\).  Each update preserves a minimizer, and nested
closed brackets whose lengths tend to zero converge to a point of the closed
minimizer set.
\end{proof}

\section{The Normal \texorpdfstring{\(\Theta\)}{Theta}: Origin, Geometry, and Computability}
\label{sec:normal-theta}

This section supplies the geometry of the motivating spectral-level-set
application.  It answers four distinct questions: what motivates a spectral
scale, which extra modeling choice creates an exact level set, why its normal is
\(\Theta=u_1v_1^\top\), and what can actually be certified when the leading
singular vectors are computed approximately.

\subsection{Spectral scaling is the motivation, not the equality constraint}

Let \(W\in\R^{m\times n}\) represent a hidden linear map \(y=Wx\).  Since
\[
 \frac{\|Wx\|_2/\sqrt m}{\|x\|_2/\sqrt n}
 \le \snorm{W}\sqrt{\frac nm},
\]
the spectral formulation of width-stable feature learning motivates the scale
\citep{yang2023spectral}
\[
 \snorm{W},\ \snorm{\Delta W}
 =\Theta\!\left(\sqrt{\frac mn}\right)
\tag{N1}
\label{eq:spectral-scaling}
\]
Equation \eqref{eq:spectral-scaling} is an asymptotic order statement.  It does
\emph{not} require the exact equality \(\snorm{W}=R\) at every iteration.
The asymptotic notation \(\Theta(\cdot)\) in
\eqref{eq:spectral-scaling} is unrelated to the matrix normal
\(\Theta(W)\) defined below; the shared glyph should not be read as a shared
mathematical object.
Following \citet{xie2026controlled}, we make the additional hard-constraint
modeling choice
\[
 \mathcal L_R
 =
 \{W\in\R^{m\times n}:\snorm{W}=R\},
 \qquad
 R\asymp\sqrt{m/n}.
\tag{N2}
\label{eq:spectral-level-set}
\]
The normal \(\Theta\) is a consequence of this additional level-set model, not
an independent optimizer ingredient and not a consequence of \(\mu\)P alone.

\subsection{The smooth stratum and its unit normal}

Write the singular values of \(W\) as
\(\sigma_1\ge\sigma_2\ge\cdots\ge\sigma_q\), where
\(q=\min\{m,n\}\).  Throughout this subsection \(R>0\).
Statements below that invoke \(\sigma_2\), a second singular subspace, or a
two-leading-value eigensolver assume \(q\ge2\).  The one-dimensional case is
elementary and is excluded from those statements rather than hidden behind an
undefined \(\sigma_2\).

\begin{lemma}[Spectral-norm differential and generalized normal]
\label{lem:spectral-normal}
If \(\sigma_1(W)>\sigma_2(W)\), then the spectral norm is real analytic in a
neighborhood of \(W\), and
\[
 D\snorm{W}[E]
 =
 \ip{u_1v_1^\top}{E}.
\tag{N3}
\label{eq:spectral-differential}
\]
Consequently
\[
 \Theta(W):=u_1(W)v_1(W)^\top
\tag{N4}
\label{eq:true-normal}
\]
is the unique outward, gradient-oriented Frobenius unit normal to the smooth
stratum of \(\mathcal L_R\); the unoriented normal space is
\(\operatorname{span}\{\Theta(W)\}\).  It is invariant under the joint sign change
\((u_1,v_1)\mapsto(-u_1,-v_1)\), has rank one, and satisfies
\[
 \snorm{\Theta}=\fnorm{\Theta}=\nnorm{\Theta}=1.
\tag{N5}
\]

If the positive leading singular value has multiplicity \(k>1\), then the
spectral norm is not differentiable and
\[
 \partial\snorm{W}
 =
 \{U_kSV_k^\top:S\succeq0,\ \tr S=1\}.
\tag{N6}
\label{eq:set-valued-weight-normal}
\]
\end{lemma}

\begin{proof}
Apply simple-eigenvalue perturbation theory to the symmetric dilation
\[
 \mathcal H(W)
 =
 \begin{pmatrix}0&W\\W^\top&0\end{pmatrix}.
\]
Its largest eigenvalue is \(\sigma_1(W)\), with normalized eigenvector
\((u_1^\top,v_1^\top)^\top/\sqrt2\).  Differentiation gives
\eqref{eq:spectral-differential}.  Simplicity makes the one-dimensional left
and right singular spaces unique up to a joint sign, which cancels in the
outer product.  The norm identities follow because \(\Theta\) has one
nonzero singular value equal to one.  Formula
\eqref{eq:set-valued-weight-normal} is the standard spectral-norm
subdifferential formula \citep{watson1992subdifferential}.
\end{proof}

\begin{proposition}[Bouligand tangent cone at a repeated leading singular value]
\label{prop:repeated-top-tangent-cone}
Let \(W\in\mathcal L_R\) have a positive leading singular value \(R\) of
multiplicity \(k\), with corresponding matrices \(U_k,V_k\).  Then the
Bouligand tangent cone of the equality level set is
\[
 T_{\mathcal L_R}^{\rm B}(W)
 =
 \left\{
 E:
 \lambda_{\max}\!\left(
 \operatorname{sym}(U_k^\top E V_k)
 \right)=0
 \right\},
\qquad
\operatorname{sym}(B)=\frac{B+B^\top}{2}.
\tag{N6a}
\label{eq:repeated-top-tangent-cone}
\]
For \(k=1\), this reduces to the hyperplane
\(\{E:\langle u_1v_1^\top,E\rangle_F=0\}\).  For \(k>1\), it is generally a
nonlinear, nonconvex cone and cannot be obtained by choosing one convenient
matrix from \(\partial\snorm{W}\).
\end{proposition}

\begin{proof}
The directional derivative of the largest singular value at a multiplicity
\(k\) cluster is
\[
 \snorm{\cdot}'(W;E)
 =
 \lambda_{\max}\!\left(
 \operatorname{sym}(U_k^\top E V_k)
 \right).
\tag{N6b}
\label{eq:spectral-directional-derivative}
\]
This follows by compressing the perturbation of the symmetric dilation
\(\mathcal H(W)\) to its leading eigenspace.

For necessity, take any \(E\in T_{\mathcal L_R}^{\rm B}(W)\).  By definition
there are \(t_j\downarrow0\) and \(E_j\to E\) such that
\(W+t_jE_j\in\mathcal L_R\).  The spectral norm is Hadamard directionally
differentiable, hence
\[
 0
 =
 \lim_{j\to\infty}
 \frac{\snorm{W+t_jE_j}-\snorm{W}}{t_j}
 =
 \snorm{\cdot}'(W;E).
\]
For sufficiency, suppose the derivative in
\eqref{eq:spectral-directional-derivative} is zero and define
\[
 W(t)=R\,\frac{W+tE}{\snorm{W+tE}}.
\]
Then \(W(t)\in\mathcal L_R\) and
\(\snorm{W+tE}=R+o(t)\), so
\(W(t)=W+tE+o(t)\).  Thus \(E\) is a Bouligand tangent direction.
\end{proof}

\begin{boundary}
\textbf{The dual/weight symmetry gap.}
The manuscript treats rank deficiency of the \emph{dual matrix}
\(G+\lambda\Theta\) set-valuedly because that inclusion is the tractable inner
problem under study.  At a repeated leading singular value of the
\emph{weight}, it does not select one element of
\(\partial\snorm{W}\) and call the resulting hyperplane the tangent space.
\Cref{prop:repeated-top-tangent-cone} shows the actual first-order object.
The present inner-oracle analysis instead excludes this stratum from any
outer claim.  That is a genuine asymmetry of scope, not a resolution of the
global stratified problem; a tangent-cone/normal-cone outer method remains
open.
\end{boundary}

\begin{remark}[A normal is not a projector]
The matrix \(\Theta\) is a unit normal, not a projector.  Under the Frobenius metric the
orthogonal tangent projector is the linear map
\[
 \Pi_{T_W}(E)=E-\ip{\Theta}{E}\Theta.
\]
\end{remark}

\begin{lemma}[Dimension-free sensitivity of the normal]
\label{lem:normal-derivative}
Let \(\gamma_W=\sigma_1(W)-\sigma_2(W)>0\).  For every perturbation \(E\),
\[
 \fnorm{D\Theta(W)[E]}
 \le
 \frac{\sqrt2}{\gamma_W}\snorm{E}.
\tag{N7}
\label{eq:normal-derivative}
\]
\end{lemma}

\begin{proof}
Choose differentiable singular-vector representatives satisfying
\(u_1^\top Du_1=v_1^\top Dv_1=0\).  The simple-eigenvector perturbation
 formula for \(\mathcal H(W)\) gives
\[
 \left\|D\left(\frac{u_1}{\sqrt2},\frac{v_1}{\sqrt2}\right)[E]\right\|_2
 \le \frac{\snorm{E}}{\gamma_W}.
\]
Since
\[
 D\Theta[E]=(Du_1[E])v_1^\top+u_1(Dv_1[E])^\top
\]
and the two summands are Frobenius orthogonal,
\[
 \fnorm{D\Theta[E]}^2
 =\|Du_1[E]\|_2^2+\|Dv_1[E]\|_2^2,
\]
which proves \eqref{eq:normal-derivative}.
\end{proof}

\subsection{Tangent constraint, ball oracle, and the role of retraction}

On the simple-leading-value stratum, the tangent space is the kernel of
\eqref{eq:spectral-differential}:
\[
 T_W\mathcal L_R
 =
 \{\Phi:\ip{\Theta(W)}{\Phi}=0\}.
\tag{N8}
\label{eq:true-tangent-space}
\]
For a step \(E=-\eta R\Phi\), this also follows from
\[
 \snorm{W-\eta R\Phi}
 =
 R-\eta R\ip{\Theta}{\Phi}
 +O\!\left(\frac{\eta^2R^2}{\gamma_W}\right).
\tag{N9}
\]
Thus the tangent LMO studied below is
\[
 p^\star
 =
 \max_{\snorm{\Phi}\le1,\ \ip{\Theta}{\Phi}=0}
 \ip{G}{\Phi}.
\tag{P}
\label{eq:normal-primal}
\]
Here \(\Theta\) is computed from the current weight \(W\), whereas \(G\) is a
loss-gradient or momentum-like matrix supplied to the oracle.  They are
different inputs.  Their only coupling in the exact dual is through the affine
line \(G+\lambda\Theta\).

\begin{lemma}[Closed ball versus unit sphere]
\label{lem:ball-sphere}
If \(p^\star>0\), every optimizer of \eqref{eq:normal-primal} saturates
\(\snorm{\Phi}=1\); hence the ball and sphere formulations have the same
optimal value and the same optimizers.  If \(p^\star=0\), the optimal values
still agree provided \(mn\ge2\), but the ball formulation also contains
interior optima.  When \(mn=1\), the tangent hyperplane is \(\{0\}\), so the
unit-sphere problem has no feasible point.
\end{lemma}

\begin{proof}
If an optimizer with positive objective had spectral norm strictly below one,
positive rescaling would remain tangent and feasible while increasing the
objective.  If \(p^\star=0\), central symmetry of the feasible set implies
that \(\ip{G}{\Phi}=0\) for every tangent feasible \(\Phi\): otherwise either
\(\Phi\) or \(-\Phi\) would have positive objective.  When \(mn\ge2\), the
nonzero normal \(\Theta\) has a nonzero Frobenius-orthogonal complement.
Choose \(E\ne0\) in that complement and replace it by
\(E/\snorm{E}\).  This gives a tangent matrix of unit spectral norm, proving
equality of the optimal values.  The \(mn=1\) exception is immediate.
\end{proof}

This lemma explains why convex duality must be established on the closed ball,
even if an implementation later normalizes a nonzero accepted direction to
unit spectral norm.

\begin{lemma}[Explicit second-order drift]
\label{lem:normal-remainder}
Let \(\gamma_W>0\) and \(\snorm{E}\le\gamma_W/4\).  Then
\[
 0
 \le
 \snorm{W+E}-\snorm{W}-\ip{\Theta}{E}
 \le
 \frac{2\snorm{E}^2}{\gamma_W}.
\tag{N10}
\label{eq:normal-remainder}
\]
In particular, a tangent step \(E=-\eta R\Phi\) with
\(\snorm{\Phi}\le1\) and \(\eta R\le\gamma_W/4\) has radial drift at most
\[
 \frac{2\eta^2R^2}{\gamma_W}.
\tag{N11}
\label{eq:linearization-drift}
\]
\end{lemma}

\begin{proof}
Convexity gives the lower bound.  In a basis whose first vector is the leading
eigenvector of \(\mathcal H(W)\), write the symmetric dilation and its
perturbation as
\[
 \mathcal H(W)=\begin{pmatrix}\sigma_1&0\\0&B\end{pmatrix},
 \qquad
 \mathcal H(E)=\begin{pmatrix}a&b^\top\\b&C\end{pmatrix},
\]
with \(\lambda_{\max}(B)\le\sigma_1-\gamma_W\) and
\(\rho:=\snorm{E}=\snorm{\mathcal H(E)}\).  The Schur-complement equation for
the perturbed leading eigenvalue gives
\[
 0\le\widetilde\sigma_1-\sigma_1-a
 \le \frac{\rho^2}{\gamma_W-2\rho}
 \le \frac{2\rho^2}{\gamma_W}.
\]
Here \(a=\ip{\Theta}{E}\).  The dilation may have additional zero eigenvalues
when \(m\ne n\); their distance from \(\sigma_1\) is at least
\(\sigma_1\ge\gamma_W\), so the same bound applies.
\end{proof}

Define the radial retraction
\[
 \Retr_W(E)=R\frac{W+E}{\snorm{W+E}}.
\tag{N12}
\label{eq:radial-retraction}
\]
Its derivative is
\[
 D\Retr_W(0)[E]
 =
 E-\frac WR\ip{\Theta}{E}.
\tag{N13}
\label{eq:retraction-derivative}
\]
Thus retraction guarantees exact radial feasibility, whereas the tangent LMO
chooses the best first-order direction under a prescribed spectral budget.
They are not interchangeable.  Retraction also does not automatically
preserve a fixed positive leading spectral gap.

\subsection{Approximate normals live in a projective geometry}

An implementation computes an approximate singular triplet
\((\hat\sigma,\hat u,\hat v)\), typically by a few power or block iterations,
and sets
\[
 \hat\Theta=\hat u\hat v^\top.
\tag{N14}
\]
For classical bilateral power iteration with a starting vector having nonzero
leading-component overlap, the right-vector angle obeys
\[
 \tan\angle(\hat v_k,v_1)
 \le
 \left(\frac{\sigma_2(W)}{\sigma_1(W)}\right)^{2k}
 \tan\angle(\hat v_0,v_1).
\tag{N14a}
\label{eq:power-rate}
\]
Thus a fixed iteration budget is reliable only when the leading ratio is
separated from one.

For a finite iteration budget, \eqref{eq:power-rate} applies with the actual
initial angle; a warm start changes that initial angle, not the exponent.
The equality \(\ip{\hat\Theta}{\Phi}=0\) is unchanged under
\(\hat\Theta\mapsto-\hat\Theta\).  The relevant object is therefore the normal
\emph{line}, not an oriented normal matrix.

Let
\[
 \alpha_u=\arccos|u_1^\top\hat u|,
 \qquad
 \alpha_v=\arccos|v_1^\top\hat v|,
 \qquad
 \alpha_u,\alpha_v\in[0,\pi/2],
\tag{N15}
\label{eq:acute-angles}
\]
and put \(\vartheta=\alpha_u+\alpha_v\).

\begin{lemma}[Exact projective distance between rank-one normals]
\label{lem:projective-normal-distance}
\[
 \min_{s\in\{\pm1\}}\nnorm{\Theta-s\hat\Theta}
 =
 2\sin\frac{\vartheta}{2},
\qquad
 \min_{s\in\{\pm1\}}\fnorm{\Theta-s\hat\Theta}
 =
 \sqrt{2-2\cos\alpha_u\cos\alpha_v}.
\tag{N16}
\label{eq:projective-distance}
\]
\end{lemma}

\begin{proof}
Choose \(s\) and a factorization
\(s\hat\Theta=\tilde u\tilde v^\top\) so that
\(u_1^\top\tilde u=\cos\alpha_u\ge0\) and
\(v_1^\top\tilde v=\cos\alpha_v\ge0\).  The difference has rank at most two.
The sum and product of its two squared singular values are
\[
 2-2\cos\alpha_u\cos\alpha_v,
 \qquad
 \sin^2\alpha_u\sin^2\alpha_v.
\]
Consequently the squared nuclear norm is
\[
 2-2\cos\alpha_u\cos\alpha_v
 +2\sin\alpha_u\sin\alpha_v
 =4\sin^2\frac{\alpha_u+\alpha_v}{2}.
\]
The other sign gives the larger of the two projective representatives.
\end{proof}

\begin{theorem}[Sharp worst-case leakage of an estimated tangent space]
\label{thm:sharp-normal-leakage}
For unit rank-one normals \(\Theta,\hat\Theta\),
\[
 \sup_{\substack{\snorm{\Phi}\le1\\
                  \ip{\hat\Theta}{\Phi}=0}}
 |\ip{\Theta}{\Phi}|
 =
 \begin{cases}
 \sin\vartheta,&0\le\vartheta\le\pi/2,\\
 1,&\pi/2\le\vartheta\le\pi.
 \end{cases}
\tag{N17}
\label{eq:sharp-normal-leakage}
\]
\end{theorem}

\begin{proof}
The absolute value permits either sign of the objective.  After the projective
orientation used in the proof of \cref{lem:projective-normal-distance}, apply
T1 strong duality (Theorem~\ref{thm:T1}) to the oracle with objective \(\Theta\) and
normal \(\hat\Theta\):
\[
 \min_{\lambda\in\R}\nnorm{\Theta+\lambda\hat\Theta}.
\]
For \(\lambda=-t\), \(t\ge0\), the squared nuclear norm is
\[
 1+t^2-2t\cos(\alpha_u+\alpha_v).
\]
Its minimum is \(\sin^2\vartheta\) when
\(\cos\vartheta\ge0\), and \(1\) otherwise.  The branch
\(\lambda=t\ge0\) has squared norm
\[
 1+t^2+2t\cos(\alpha_u-\alpha_v),
\]
whose minimum is \(1\) because
\(|\alpha_u-\alpha_v|\le\pi/2\).  Comparing the branches proves
\eqref{eq:sharp-normal-leakage}.
\end{proof}

\begin{boundary}
\Cref{thm:sharp-normal-leakage} is a sharp \emph{uniform worst-case upper
bound}.  It is not a lower bound on every realized oracle direction: a
particular direction may be tangent to both normal hyperplanes.  Likewise, an
angle measured directly between two normal matrices is not, in general, the
sum \(\alpha_u+\alpha_v\) appearing in \eqref{eq:sharp-normal-leakage}.  No
empirical leakage percentage is inferred in this paper because the required
pair \((\alpha_u,\alpha_v)\) was not recorded.
\end{boundary}

\section{A Feasible Certificate and Its Finite-Precision Boundary}
\label{sec:certificates}

An approximate matrix output cannot be certified merely by reporting a polar
residual and a tangent residual.  If the tangent residual is nonzero, the
matrix is not feasible for the primal oracle, so its objective value need not
be a valid primal lower bound.

Assume here that \(\Theta=uv^\top\) is a unit rank-one normal.  Given an
arbitrary approximate matrix \(Z\), define the spectral normalization
\[
\bar Z
=
\frac{Z}{\max\{1,\snorm{Z}\}},
\qquad
c=\ip{\Theta}{\bar Z},
\tag{54}
\label{eq:Zbar-c}
\]
and the feasible repair
\[
\widetilde Z
=
\frac{\bar Z-c\Theta}{1+|c|}.
\tag{55}
\label{eq:Ztilde}
\]

\begin{theorem}[T6: feasible repair and valid primal--dual gap in exact arithmetic]
\label{thm:T6}
The repaired matrix satisfies
\[
\ip{\Theta}{\widetilde Z}=0,
\qquad
\snorm{\widetilde Z}\le1.
\tag{56}
\label{eq:T6-feas}
\]
For any multiplier \(\lambda\), define
\[
\Gap(\lambda,Z)
=
\nnorm{G+\lambda\Theta}
-\ip{G}{\widetilde Z}.
\tag{57}
\label{eq:true-gap}
\]
Then
\[
\Gap(\lambda,Z)
\ge
p^\star-\ip{G}{\widetilde Z}
\ge0.
\tag{58}
\label{eq:gap-cover}
\]
Thus \eqref{eq:true-gap} is a valid upper bound on the primal LMO
suboptimality.

If
\[
r_{\mathrm{pol}}
=
\nnorm{G+\lambda\Theta}
-\ip{G+\lambda\Theta}{\bar Z},
\tag{59}
\label{eq:rpol}
\]
then \(r_{\mathrm{pol}}\ge0\), and
\[
\Gap(\lambda,Z)
\le
r_{\mathrm{pol}}
+|\lambda||c|
+\frac{2\nnorm{G}|c|}{1+|c|}.
\tag{60}
\label{eq:gap-decomposition}
\]
Finally, if a computable quantity \(\widehat\phi(\lambda)\) has been proved to
satisfy
\[
\widehat\phi(\lambda)\ge\nnorm{G+\lambda\Theta},
\]
then
\[
\widehat\phi(\lambda)-\ip{G}{\widetilde Z}
\tag{61}
\label{eq:upper-certificate}
\]
is also a valid certificate.
\end{theorem}

\begin{proofidea}
First normalize into the spectral ball.  Then subtract the normal component
and divide by a worst-case triangle-inequality factor so that tangency and
spectral feasibility both hold exactly.  Only after this repair is the
dual-minus-primal difference a legitimate gap.
\end{proofidea}

\begin{intuition}
T6 turns any approximate polar map---including one affected by truncation,
finite precision, or an imperfect multiplier---into a direction that the
downstream method may safely use.  The certificate measures the price of the
repair rather than silently treating an infeasible matrix as primal feasible.
\end{intuition}

\begin{proof}
By definition, \(\snorm{\bar Z}\le1\).  Because the rank-one unit normal
satisfies \(\ip{\Theta}{\Theta}=\fnorm{\Theta}^2=1\),
\[
\ip{\Theta}{\widetilde Z}
=
\frac{c-c}{1+|c|}
=0.
\]
Furthermore,
\[
\snorm{\bar Z-c\Theta}
\le
\snorm{\bar Z}+|c|\snorm{\Theta}
\le1+|c|.
\]
Dividing by \(1+|c|\) proves \eqref{eq:T6-feas}.

Weak duality from Theorem~\ref{thm:T1} gives
\[
\nnorm{G+\lambda\Theta}\ge p^\star.
\]
Since \(\widetilde Z\) is primal feasible,
\[
\ip{G}{\widetilde Z}\le p^\star.
\]
Subtracting proves \eqref{eq:gap-cover}.

For the decomposition, put \(A=G+\lambda\Theta\).  Adding and subtracting
\(\ip{A}{\bar Z}\) gives
\[
\begin{aligned}
\Gap(\lambda,Z)
&=
\nnorm{A}-\ip{G}{\widetilde Z}\\
&=
r_{\mathrm{pol}}
+\lambda c
+\ip{G}{\bar Z-\widetilde Z}.
\end{aligned}
\]
Equation \eqref{eq:Ztilde} gives
\[
\bar Z-\widetilde Z
=
\frac{|c|\bar Z+c\Theta}{1+|c|}.
\]
By spectral/nuclear duality,
\[
\begin{aligned}
\left|\ip{G}{\bar Z-\widetilde Z}\right|
&\le
\frac{|c|\nnorm{G}\snorm{\bar Z}
      +|c|\nnorm{G}\snorm{\Theta}}
     {1+|c|}\\
&\le
\frac{2\nnorm{G}|c|}{1+|c|}.
\end{aligned}
\]
Using \(\lambda c\le|\lambda||c|\) proves
\eqref{eq:gap-decomposition}.  Replacing the exact dual value by any proved
upper bound preserves the coverage inequality, giving
\eqref{eq:upper-certificate}.
\end{proof}

\begin{boundary}
The feasibility and coverage statements above are exact-real-arithmetic
identities.  A machine-rigorous certificate additionally needs outward-rounded
upper bounds for the nuclear norm and spectral normalization, a downward-
rounded primal objective, and enclosures for the tangency calculation.  Plain
double-precision evaluation is a numerical check, not by itself a validated
certificate.  These are necessary ingredients, not a complete validated
implementation: an interval enclosure of the tangency residual alone does
not certify exact membership in the hyperplane unless the final repair is
itself performed exactly or with an interval-safe correction.
\end{boundary}

\begin{corollary}[Optional unit-spectral-norm output]
\label{cor:unit-repair}
Let
\[
P=\bar Z-c\Theta.
\]
If \(P\ne0\), choose
\[
s_G=
\begin{cases}
1,&\ip{G}{P}\ge0,\\
-1,&\ip{G}{P}<0,
\end{cases}
\qquad
Z_{\rm unit}=s_G\frac{P}{\snorm{P}}.
\tag{62}
\label{eq:unit-repair}
\]
Then
\[
\ip{\Theta}{Z_{\rm unit}}=0,\qquad
\snorm{Z_{\rm unit}}=1,
\]
and
\[
\ip{G}{Z_{\rm unit}}\ge\ip{G}{\widetilde Z}.
\tag{63}
\label{eq:unit-repair-objective}
\]
Consequently, replacing \(\widetilde Z\) by \(Z_{\rm unit}\) never worsens the
valid primal--dual certificate.

Likewise, if \(p^\star>0\), then the smooth solution
\(\Phi_\delta\ne0\) may be replaced by
\[
\Phi_\delta^{\rm unit}
=\frac{\Phi_\delta}{\snorm{\Phi_\delta}}.
\]
It remains tangent, has unit spectral norm, and has no smaller original
linear objective.  If
\[
\sigma_1(G+\lambda_\delta\Theta)\ge s_0>0,
\]
then the rescaling factor satisfies
\[
\frac{1}{\snorm{\Phi_\delta}}
=
\sqrt{1+\frac{\delta^2}
{\sigma_1(G+\lambda_\delta\Theta)^2}}
=1+O(\delta^2).
\tag{64}
\label{eq:smooth-unit-factor}
\]
\end{corollary}

\begin{proof}
Tangency and unit norm in \eqref{eq:unit-repair} are immediate.  Moreover,
\(\snorm{P}\le1+|c|\), so if \(\ip{G}{P}\ge0\),
\[
\ip{G}{Z_{\rm unit}}
=\frac{\ip{G}{P}}{\snorm{P}}
\ge\frac{\ip{G}{P}}{1+|c|}
=\ip{G}{\widetilde Z}.
\]
If \(\ip{G}{P}<0\), the left side is positive whereas the right side is
negative.

For the smooth direction, comparison with the feasible zero direction in
\eqref{eq:regularized-primal} gives
\[
\ip{G}{\Phi_\delta}
\ge\delta R(\Phi_\delta)\ge0.
\]
To see rigorously that \(\Phi_\delta\ne0\) when \(p^\star>0\), let
\(\Phi^\star\) be an exact primal optimizer.  For sufficiently small
\(t>0\), the exact identity
\(1-\sqrt{1-z^2}=z^2/(1+\sqrt{1-z^2})\le z^2\) gives
\[
 R(t\Phi^\star)\le t^2\|\Phi^\star\|_F^2.
\]
Consequently
\[
 \ip{G}{t\Phi^\star}-\delta R(t\Phi^\star)
 \ge tp^\star-\delta t^2\|\Phi^\star\|_F^2>0,
\]
whereas the regularized objective at zero is zero.  Hence the unique smooth
optimizer is nonzero, and positive radial rescaling cannot decrease its
linear objective.  Finally,
\eqref{eq:Zdelta} gives
\[
\snorm{\Phi_\delta}
=
\frac{\sigma_1(G+\lambda_\delta\Theta)}
{\sqrt{\sigma_1(G+\lambda_\delta\Theta)^2+\delta^2}},
\]
which proves \eqref{eq:smooth-unit-factor}.
\end{proof}

\begin{boundary}
The unit repair is certified primal post-processing and an experimental
scale-control device.  The map \(P\mapsto P/\snorm{P}\) is discontinuous at
\(P=0\) and may amplify numerical noise nearby.  Moreover,
\(\Phi_\delta^{\rm unit}\) is no longer the exact optimizer of the
regularized primal at finite \(\delta\), so its Fenchel center-path identity
must not be asserted pointwise.  Under the lower bound in
\eqref{eq:smooth-unit-factor}, however, its rescaling is \(1+O(\delta^2)\)
and it has the same limiting center.
\end{boundary}

\subsection{How to read the decomposition}

The three terms in \eqref{eq:gap-decomposition} have different meanings.

\begin{enumerate}
  \item \(r_{\mathrm{pol}}\) measures matrix-function or polar optimality
        error.
  \item \(|\lambda||c|\) measures the coupling between multiplier size and
        tangent residual.
  \item The last term is the explicit price of moving the infeasible output
        into the primal feasible set.
\end{enumerate}

\paragraph{Certified stopping rule.}
Given a target inner error budget \(\varepsilon_t\), stop the inner oracle only
when a valid quantity of the form \eqref{eq:true-gap} or
\eqref{eq:upper-certificate} is at most \(\varepsilon_t\).  Reporting only
\(|c|\) is not sufficient.

\subsection{A certifiable reference homotopy oracle}

The theory suggests the following implementation pattern.

\begin{enumerate}
  \item \textbf{Choose a smoothing level.}
        Start from a moderate \(\delta\); warm-start it from the previous
        oracle call when appropriate.
  \item \textbf{Bracket the smooth multiplier.}
        Use the proved multiplier bound
        \eqref{eq:smooth-lambda-bound}, or a tighter bracket whose coverage has
        been verified.
  \item \textbf{Solve the continuous root.}
        Apply bisection or a safeguarded scalar solver to
        \(h_\delta(\lambda)=0\).  Unlike the compact selection, this function
        is continuous and strictly increasing.
  \item \textbf{Form the smooth direction.}
        Evaluate \(Z_\delta(G+\lambda\Theta)\).  If the scalar root is only
        approximate or the matrix function is computed in finite precision,
        apply \eqref{eq:Zbar-c}--\eqref{eq:Ztilde}.
  \item \textbf{Certify.}
        Compute \eqref{eq:true-gap}, or a proved upper version
        \eqref{eq:upper-certificate}.
  \item \textbf{Refine only when needed.}
        If the certificate exceeds the requested tolerance \(\varepsilon_t\), reduce
        \(\delta\), tighten the scalar root, or switch to the exact
        set-valued SVD oracle.  Otherwise accept the repaired direction.
\end{enumerate}

This homotopy is ``accuracy on demand'' only after the normal accuracy and any
assumptions required by a downstream method have been certified independently.
T6 makes only the \emph{inner} oracle gap computable; it does not certify the
normal or an outer descent model.
Consequently this is a certifiable reference workflow, not a claim of a
fully deployable per-step large-model routine.

\begin{remark}[Conservatism]
The repair is deliberately simple and globally safe.  It may be conservative:
projecting onto the exact intersection of a spectral ball and a hyperplane
could produce a tighter primal value, but would itself require a nontrivial
optimization.  T6 prioritizes a closed-form guarantee.
\end{remark}

\section{Detailed numerical verification}
\label{sec:supp-numerics}

The local laws of Section~\ref{sec:rates} predict constants, not merely exponents, so
they can be falsified numerically.  We report closed-form instances in which
the displayed constants are available in exact arithmetic.  The
\(2\times2\) checks evaluate the smoothed scalar root of
\eqref{eq:smooth-root-main} in high-precision decimal arithmetic through the
closed principal-square-root formula for
\(Z_\delta(A)=A(A^\top A+\delta^2I)^{-1/2}\); the diagonal corank-two
instance is evaluated entrywise.  The roots are located by
bisection on the strictly increasing function \(h_\delta\); assertion-based
scripts accompany the paper.

For the first group, \(\Theta=e_1e_1^\top\) and
\(G=\bigl(\begin{smallmatrix}0&b_1\\b_2&c\end{smallmatrix}\bigr)\).

\begin{center}
\begin{tabular}{@{}llll@{}}
\toprule
Regime & Instance and predicted constant & Smoothing & Computed \\
\midrule
Theorem~\ref{thm:rate-regular} &
\(b_1=b_2=1,\ c=\tfrac12\);\ \(\lambda^\star=\tfrac12\), \(\kappa=\tfrac12\) &
\(\delta=10^{-4}\) & \(-1.777777665\)\\
\(O(\delta^2)\) &
\((\lambda_\delta-\lambda^\star)/\delta^2\to-b^\star/\kappa=-\tfrac{16}{9}\)
&\(\delta=10^{-5}\) & \(-1.777777777\)\\
\addlinespace
Corollary~\ref{thm:rate-strict} &
\(b_1=b_2=\tfrac12,\ c=1\);\ \(\lambda^\star=\tfrac14\),
\(a=\tfrac15\), \(\beta=\tfrac45\) &
\(\delta=10^{-6}\) & \(-0.3227484177\)\\
\(\Theta(\delta)\) &
\((\lambda_\delta-\lambda^\star)/\delta\to t_0=-\tfrac{5}{4\sqrt{15}}\)
& \(\delta=10^{-8}\) & \(-0.3227486102\)\\
\addlinespace
Theorem~\ref{thm:rate-boundary} &
\(b_1=b_2=c=1\);\ \(\lambda^\star=1\), \(a=\beta=\tfrac12\), \(\kappa=\tfrac12\) &
\(\delta=10^{-9}\) & \(1.259920256\)\\
\(\Theta(\delta^{2/3})\) &
\((\lambda^\star-\lambda_\delta)/\delta^{2/3}\to(2\beta\kappa)^{-1/3}=2^{1/3}\)
& \(\delta=10^{-12}\) & \(1.259921042\)\\
\bottomrule
\end{tabular}

\smallskip
{\small
Exact values: \(-16/9=-1.777777778\),
\(-5/(4\sqrt{15})=-0.3227486122\),
\(2^{1/3}=1.259921050\).}
\end{center}

For the same boundary data, standard Huber gives
\((\lambda^\star-\lambda_\delta^{\rm hub})/\delta=1.9999999656\) at
\(\delta=10^{-8}\), versus \(1/\beta=2\).  The scaled-Huber regression in
the code also confirms that moving the first saturation threshold to
\(L=1/2\) changes the constant to \(L/\beta=1\), as required by
\eqref{eq:saturated-boundary-rate-main}.

The preceding instances are \(2\times2\).  To exercise the general
completion and arbitrary-corank strict law, take
\[
 A^\star=\diag(1,1,0,0),\qquad
 \Theta=\diag\!\left(\tfrac38,\tfrac38,1,\tfrac12\right),\qquad
 \lambda^\star=1,\qquad G=A^\star-\Theta.
\]
Then \(a=3/4\), \(C=\diag(1,1/2)\), and \(\beta=3/2\).  Hard
water-filling gives
\(\eta_F=3/5\),
\(K_F=-\diag(3/5,3/10)\), and
\(\|K_F\|_F=\sqrt{0.45}\).  The smooth-center equation gives
\[
 \eta_R=0.715080068980,\qquad
 K_R=-\diag(0.581665969555,0.336668060890),
\]
with \(\|K_F-K_R\|_F=0.040996138377\).  Writing
\(\Phi_F^\star\) and \(\Phi_R^\star\) for the two corresponding completed
certificates, the actual smoothing path gives
\begin{center}
\begin{tabular}{@{}llll@{}}
\toprule
\(\delta\) & \(\lambda_\delta-\lambda^\star\) &
\(\|\Phi_\delta-\Phi_R^\star\|_F\) &
\(\|\Phi_\delta-\Phi_F^\star\|_F\)\\
\midrule
\(10^{-2}\) & \(-7.150296\!\times\!10^{-3}\) & \(7.896\!\times\!10^{-5}\) & \(4.0990\!\times\!10^{-2}\)\\
\(10^{-3}\) & \(-7.150796\!\times\!10^{-4}\) & \(7.859\!\times\!10^{-7}\) & \(4.0996\!\times\!10^{-2}\)\\
\(10^{-4}\) & \(-7.150801\!\times\!10^{-5}\) & \(7.855\!\times\!10^{-9}\) & \(4.0996\!\times\!10^{-2}\)\\
\bottomrule
\end{tabular}
\end{center}
Thus the path converges to the smooth rather than the minimum-Frobenius
center, while
\((\lambda_\delta-\lambda^\star)/\delta\to-\eta_R\), as predicted by
Theorems~\ref{thm:general-fenchel-center} and~\ref{thm:rate-strict-general}
and Corollary~\ref{cor:center-agreement-main}.
On the same line, the standard Huber path gives slope
\(-\eta_F=-3/5\) and converges to \(\Phi_F^\star\) to below
\(10^{-11}\) at \(\delta=10^{-5}\), verifying
Proposition~\ref{prop:smoothing-family-main} on the same nonsingleton slice.

The preceding corank-two line is diagonal: the relevant singular dyads are
fixed.  In the notation of Proposition~\ref{prop:strict-constants-closed-main},
both \(P_{\rm big}(s)\) and \(\widetilde B(s)\) are constant.  Hence
\(\partial_s\mathcal W=0\), which proves \(t_1=0\) and
\(\mathcal M_R=0\) rather than merely observing their cancellation.  To test a nonzero first-order
coefficient without random data, take
\[
 A^\star=\diag(1,3/5,0,0),\qquad \lambda^\star=1,
\]
\[
 \Theta=
 \begin{pmatrix}
  1/6&0&2/5&0\\
  0&1/6&0&3/10\\
  1/5&0&2/3&0\\
  0&1/10&0&1/3
 \end{pmatrix},
 \qquad G=A^\star-\Theta.
\]
In the displayed coordinates, \(a=1/3\),
\(C=\diag(2/3,1/3)\), \(\beta=1\), and
\(\eta_R=0.644812526657\).  At \(\delta=10^{-5}\), the extended-precision
multiplier ratio is \(-0.644805516183\), while the deterministic direction
ratio is \(0.671768254\),
respectively, already exhibiting a nonzero linear direction term.  Direct evaluation of
\eqref{eq:t1-closed-main}--\eqref{eq:MR-closed-main} gives
\[
 t_1=0.7010530600,
 \qquad \|\mathcal M_R\|_F=0.6717753807.
\]
At \(\delta=10^{-6}\), the observed direction ratio is
\(0.6717760680\), and the cosine between
\((\Phi_\delta-\Phi_R^\star)/\delta\) and \(\mathcal M_R\) equals one to
ten displayed digits.  The corresponding 50-digit finite-scale estimate of
\(t_1\) at \(\delta=10^{-5}\) is \(0.7010473954\).  Thus the table verifies the
predicted vector coefficient, not only its exponent.  This demonstrates
sharpness for one fixed instance; no generic measure statement is inferred.

The compact-root branch of
Corollary~\ref{cor:compact-root-superconvergence-main} is sharp as well.
For \(A^\star=\diag(2,1,0)\), \(\Theta=\diag(1,-1,1)\), and
\(G=A^\star-\Theta\), one has \(a=0\), \(C=(1)\),
\(J_\star=\diag(1/4,1,0)\), and \(\chi_\star=-3/4\).  Thus the predicted
coefficients are \(-3/8\) and
\(\|\mathcal N_0\|_F=\sqrt{26}/8=0.6373774392\ldots\).
At \(\delta=10^{-4}\), high-precision bisection gives the normalized pair
\((-0.3749999965,\,0.6373774341)\); the assertion script checks monotone
convergence over four scales.  Thus, even when the compact-polar equation
already has a root, the nearby smooth path has nonzero errors beginning at
cubic and quadratic order.

The doubly cancelled branch in
\eqref{eq:compact-root-quintic-main} uses
\[
 A^\star=\diag(3,2,1,0),\qquad \lambda^\star=1,
\]
\[
 \Theta=\begin{pmatrix}
  27/10&0&0&3/10\\
  0&-16/5&0&1/10\\
  0&0&1/2&1/4\\
  1/5&-3/20&1/10&1
 \end{pmatrix},\qquad G=A^\star-\Theta.
\]
Here \(a=0\), \(C=(1)\), \(d_0=1\),
\(\chi_\star=0\), and \(\chi_2=1/3\).  Hence the predicted quintic
coefficient is \(-1/8\), and the quartic direction coefficient after
subtracting \(-\delta^2J_\star/2\) is
\[
 \diag\!\left(\frac1{216},\frac3{128},\frac38,-\frac18\right).
\]
Extended-precision bisection gives
\begin{center}
\begin{tabular}{@{}lllll@{}}
\toprule
\(\delta\) & \(10^{-2}\) & \(10^{-3}\) & \(10^{-4}\) & \(10^{-5}\)\\
\midrule
\((\lambda_\delta-\lambda^\star)/\delta^5\) &
\(-0.1248207361\)&\(-0.1249833498\)&\(-0.1249983476\)&\(-0.1249998349\)\\
\bottomrule
\end{tabular}
\end{center}
At \(\delta=10^{-4}\), the diagonal of the normalized quartic direction is
\[
 (0.004629629625,,0.023437499951,,0.374999996875,
   -0.124998347562),
\]
while the Frobenius norm of its off-diagonal part is below
\(4\times10^{-6}\).  The high-precision assertion script checks all four
scales and the full matrix coefficient.

For Theorem~\ref{thm:crossover-main} we use the one-parameter family
\(b_1=b_2=1\), \(c=1+\varepsilon\), for which
\(\beta_0=\kappa_0=\tfrac12\) and
\(\Delta_\varepsilon=(c^2-1)/(1+c^2)\).  Choosing \(\varepsilon=
\varepsilon(\delta)\) so that \(\Delta_\delta=\zeta\delta^{2/3}\) places the
family in the critical window \eqref{eq:crossover-critical-main}, whose
limit solves \(\tfrac12y^3+\zeta y^2=1\).

\begin{center}
\begin{tabular}{@{}lllll@{}}
\toprule
\(\zeta\) & \(0\) & \(1\) & \(3\) & \(10\)\\
\midrule
predicted \(y_\zeta\) &
\(1.259921050\)&\(0.839286755\)&\(0.552474611\)&\(0.313775961\)\\
computed \(x_\delta/\delta^{2/3}\), \(\delta=10^{-9}\) &
\(1.259920256\)&\(0.839284833\)&\(0.552471210\)&\(0.313769661\)\\
\bottomrule
\end{tabular}
\end{center}

For a valid strict-side moving path, take
\(\Delta_\delta=\delta^{1/3}\); then
\(\Delta_\delta\to0\) while
\(\Delta_\delta/\delta^{2/3}\to\infty\).  The normalization in
\eqref{eq:crossover-inner-main}, whose limit is one for this family, gives
\begin{center}
\begin{tabular}{@{}lllll@{}}
\toprule
\(\delta\) & \(10^{-6}\) & \(10^{-9}\) & \(10^{-12}\) & \(10^{-15}\)\\
\midrule
\(x_\delta\sqrt{\Delta_\delta}/\delta\) &
\(0.979950782\)&\(0.997994100\)&\(0.999799770\)&\(0.999979992\)\\
\bottomrule
\end{tabular}
\end{center}
This path respects the moving-family quantifier.  By contrast, fixing
\(\varepsilon\) and sending \(\delta\) to zero tests the fixed strict-kink
law, not the crossover theorem.

Finally, the oracle returns a matrix direction, not only a multiplier.  At
the boundary instance \(b_1=b_2=c=1\), direct differentiation gives
\[
 \|\mathcal M_0\|_F=\sqrt{3/2},
 \qquad
 \frac{\|\Phi_\delta-\Phi^\star\|_F}{x_\delta}
 \longrightarrow\sqrt{3/2}.
\]
At \(\delta=10^{-15}\), the computed ratio is
\(1.224744871327\), versus
\(\sqrt{3/2}=1.224744871392\).  The same verification includes critical,
strict-side, boundary, and deliberately oscillatory paths for the unified
ratio \(x_\delta/\widehat x_\delta\).

We emphasize what these tables do and do not support.  They confirm the
exponents and leading constants of the local laws on instances where those
constants are known in closed form.  They are not evidence about how often
rank-deficient crossings occur in any application, and no such claim is made.


{\footnotesize
\setlength{\bibsep}{0pt plus 0.2ex}
\bibliographystyle{plainnat}
\bibliography{references}
}

\end{document}
